\section*{Ethics Considerations}

This work studies an attack on safety-critical learning-enabled control, so we
state its scope and safeguards explicitly.  All experiments run in simulation
(Safe-Control-Gym with PyBullet dynamics); no physical plant, vehicle, or
deployed industrial system was involved, and no human subjects or personal data
were used.  The reported ``failures'' are simulated constraint violations used
to compare deployment contracts, not incidents on any real system.

The evaluated threat begins \emph{after} a reward-ingestion principal in the
asynchronous learning-data plane has been compromised, and grants the attacker
only bounded scalar writes to reward records.  We deliberately do not
demonstrate, and do not provide, any capability to obtain that foothold: we
present no exploit against a historian, broker, replay service, or reward
producer, and no technique to evade log-integrity controls.  The contribution
is thus not a new intrusion capability but a characterization of what a
\emph{known} least-privilege data-integrity gap implies for certificate
inheritance at raw-policy release.

We judge the disclosure to be net beneficial.  The same study that exposes the
assurance-contract gap also identifies concrete, deployable countermeasures and
measures their cost: origin-bound append-only reward records and trusted
recomputation remove the upstream precondition entirely (recomputation blocks
every evaluated edit), while resident predictive authority, commit admission,
and permanent runtime filtering contain the downstream physical consequence.
The effective attack is also readily detectable by simple task-aware reward
checks.  Surfacing this contract distinction, together with its defenses, before
learning-enabled controllers are widely promoted from protected adaptation to
unmonitored release reduces rather than increases real-world risk.  We are not
aware of a specific fielded system exhibiting this exact pipeline, so no vendor
disclosure was applicable; operators adopting protected-adaptation-then-release
workflows are the intended audience for the mitigations.

\section*{Open Science}

We support the NDSS open-science policy.  The artifact includes row-level result
CSVs, the exact commands and seed namespaces to reproduce every table and
figure, upstream revisions and model files, dependency versions, source-linked
generated tables, and SHA-256 checksums for all locked files, together with an
automated audit that checks table provenance, claim locks, and format.  We
intend to release the code and data required to reproduce the paper's results
under an open license upon publication.
